% 《大学物理简明教程》第5版第9章填空题9.2.1—9.2.14
% 原书/PDF第238—239页；据用户扫描页校订题意并独立推导。

\subsection{填空题：感应电动势、位移电流及磁场能量}

\sourceproblem{9.2.1}{238}{238}
\problemtext 把一金属圆环从磁极间沿垂直于磁感应强度的方向抽出，
问圆环会受到哪种电磁作用。
\workings
闭合环路在移出磁场边界时，穿过回路的磁通量随时间变化，
$\mathcal E=-d\Phi_B/dt$，由此产生感应电流。
感应电流在磁场内受安培力，按楞次定律该合力阻碍移出运动。
因此填空为\textbf{磁阻力（阻碍抽出的安培力）}。
\crosscheck 只有圆环\emph{部分跨越磁场边界}才有非零磁通量变化；
若整环在空间均匀、无限延伸的磁场中匀速平移，则总电动势为零。

\sourceproblem{9.2.2}{238}{238}
\problemtext 分别填出动生电动势、感生电动势对应的非静电场力来源，
以及产生感生电场的场源。
\workings
动生电动势来自运动电荷所受磁力
$\boxed{\bm F_{\mathrm{mot}}=q\bm v\times\bm B}$；
感生电动势来自涡旋电场对电荷的力
$\boxed{\bm F_{\mathrm{ind}}=q\bm E_{\mathrm{ind}}}$；
感生电场的场源是\textbf{随时间变化的磁场}，满足
\[
\boxed{\nabla\times\bm E_{\mathrm{ind}}
=-\frac{\partial\bm B}{\partial t}}.
\]
\crosscheck 电磁感应有动生和感生两类物理机制，
不可全部写作洛伦兹磁力或全部写作静电库仑力。

\sourceproblem{9.2.3}{238}{238}
\problemtext 长$\ell$金属直导线在垂直匀强磁场$\bm B$的平面内
绕固定轴以角速度$\omega$转动。
转轴取在导线的哪里时两端电动势最大，最大值是多少；
取在何处两端电动势最小，最小值是多少？
\workings
设转轴距导线一端为$a$，
从轴向该端和另一端分别有距离$a,\ell-a$，
有向端电势差
\[
\mathcal E=\frac12B\omega
\bigl[(\ell-a)^2-a^2\bigr]
=\frac12B\omega(\ell^2-2\ell a).
\]
在$a=0$或$a=\ell$时，
\[
\boxed{|\mathcal E|_{\max}=\frac12B\omega\ell^2},
\]
转轴位于\textbf{导线端点}；
当$a=\ell/2$即转轴位于\textbf{中点}时，
\[
\boxed{|\mathcal E|_{\min}=0}.
\]
\crosscheck 两半段虽各自具有对轴的动生电动势，
但关于轴的有向贡献相反，故绕中点转动时\emph{端到端}电动势零。

\sourceproblem{9.2.4}{238}{238}
\problemtext 通过单匝线圈的磁通量为
$\Phi(t)=6t^2+9t+2$（SI单位），求$t=2$ s的感应电动势。
\workings
法拉第定律给
\[
\mathcal E(t)=-\frac{d\Phi}{dt}=-(12t+9)\,\mathrm V,
\quad\boxed{\mathcal E(2\,\mathrm s)=-33\,\mathrm V}.
\]
若只求电动势\emph{大小}则是$\boxed{33\,\mathrm V}$。
\crosscheck 负号是相对于预设线圈正绕行方向的定向符号；
原题没有给明确的绕行方向时应同时说明大小和有向取值约定。

\sourceproblem{9.2.5}{238}{238}
\problemtext 写出法拉第电磁感应定律的闭合积分表达式。
\workings
静止闭合路径的场方程为
\[
\boxed{\oint_{\partial S}\bm E\cdot d\bm l
=-\frac d{dt}\int_S\bm B\cdot d\bm S
=-\frac{d\Phi_B}{dt}}.
\]
\crosscheck 环路和面的正向应满足右手关系；
若边界随时间运动，还必须将运动边界的
$\bm v\times\bm B$项计入总电动势。

\sourceproblem{9.2.6}{238}{238}
\problemtext 写出时间变化磁场与它所激发的涡旋电场的微分关系。
\workings
麦克斯韦--法拉第方程为
\[
\boxed{\nabla\times\bm E=-\frac{\partial\bm B}{\partial t}}.
\]
\crosscheck 左边单位为$\mathrm{V/m^2}$，
右边单位为$\mathrm{T/s}$，两者在SI制下等价；
这也说明一般感生电场不满足无旋静电场条件。

\sourceproblem{9.2.7}{238}{238}
\problemtext 无限长直导线竖直放置、电流向上；
在其右侧有同平面的任意形状闭合线圈。
(1) 线圈整体沿导线向下运动；
(2) 线圈在本平面内以垂直导线的速度向导线靠近。
约定\emph{顺时针}感应电动势为正，分别判定正负。
\workings
(1) 无限长导线电场与磁场在竖直方向平移不变，
线圈平行移动时磁通量不变：
\[
\boxed{\mathcal E_1=0}.
\]
(2) 右侧磁场指向纸内，
线圈靠近导线使进入纸面的通量\emph{增大}，
感应电流应产生指向纸外的反向场，故沿\emph{逆时针}。
题设顺时针为正，所以
\[
\boxed{\mathcal E_2<0}.
\]
\crosscheck 特别区分约定：若规定逆时针为正，
第二项的\emph{数值符号}会反转，
但实际绕行方向不会改变。

\sourceproblem{9.2.8}{238}{238}
\problemtext 一$N$匝矩形线圈固定于小振动杆，
其宽度（垂直运动方向）为$b$，有一部分位于恒定垂直磁场$B$内。
杆的位置为$x(t)=A\cos\omega t$，
求忽略边缘效应时线圈中的感应电动势。
\workings
以边缘内的重叠长度$x$表示有效面积，
可取$\Phi_B=NBb\,x$（若图中$x$为向外位移则相应变号）。
因此
\[
\boxed{\mathcal E(t)=-NBb\dot x
=NBb\,A\omega\sin\omega t}.
\]
最大大小为$\boxed{\mathcal E_{\max}=NBbA\omega}$。
\crosscheck 这一式要求线圈的有源边一直跨越均匀磁场的\emph{边界}，
且另一边不同时跨界；否则有效面积函数不是简单的$b\,x$。

\sourceproblem{9.2.9}{238}{238}
\problemtext 在两条平行金属导轨上有可自由滑动的导体杆
$AB$（左）及$CD$（右），穿过纸面的均匀磁场向里，
外力使左杆$AB$向左运动。
问右杆$CD$将如何运动、磁场对左杆施加的合力指向何处。
\workings
$AB$向左移动使两杆组成的回路面积\emph{增加}，
向纸内的磁通量增加。
楞次定律要求感应电流产生向纸外的磁场，
因此整个回路按逆时针方向通电：
左杆中电流由$A$流向$B$（向下），
右杆$CD$由$D$流向$C$（向上）。
由$\bm F=I\bm l\times\bm B$，
右杆受力向\textbf{左}，左杆受力向\textbf{右}。
\crosscheck 左杆受磁力总与外力驱动的左移方向相反，
符合感应电流消耗机械输入能量的楞次定律。

\sourceproblem{9.2.10}{238}{238}
\problemtext 半径$r=20$ cm的圆形闭合回路位于垂直磁场
$B=1.0$ T中，半径从$t=0$以恒定速率
$dr/dt=-50$ cm/s缩小。
求初始感应电动势大小，以及让电动势保持该值时
需要的恒定面积变化率。
\workings
磁通量$\Phi=B\pi r^2$，
故
\[
|\mathcal E|=B\left|2\pi r\frac{dr}{dt}\right|
=1\cdot2\pi(0.20)(0.50)
=\boxed{0.20\pi\,\mathrm V\approx0.628\,\mathrm V}.
\]
若$B$不变，则要求
\[
\boxed{\frac{dS}{dt}=-\frac{|\mathcal E|}{B}
=-0.20\pi\,\mathrm{m^2/s}}.
\]
\crosscheck 保持电动势恒定意味着$S(t)$线性减小；
半径的收缩速率需随$r$减小而增大，
不能一直保持$dr/dt=-0.5$ m/s。

\sourceproblem{9.2.11}{238}{238}
\problemtext 细长螺线管横截面积$S=1\,\mathrm{cm^2}$，
总匝数$N=400$，电流$I=4$ A时内部均匀场$B=0.2$ T。
忽略漏磁，求自感。
\workings
总磁链$\Psi=NB S$，
自感定义$L=\Psi/I$，因此
\[
\boxed{L=\frac{NBS}{I}
=\frac{400\cdot0.2\cdot10^{-4}}4
=2.0\times10^{-3}\,\mathrm H}.
\]
\crosscheck $1\,\mathrm{cm^2}=10^{-4}\,\mathrm{m^2}$；
此处$B$已给，不必另外知道螺线管长度。

\sourceproblem{9.2.12}{238}{238}
\problemtext 任意两个固定几何形状的线性互感回路1、2，
流过相同的稳恒电流$I$。
由回路2电流穿过回路1的磁通量记$\Phi_{12}$，
反之记$\Phi_{21}$。求二者关系。
\workings
互感互易性给
$M_{12}=M_{21}=M$，
若每个回路都是单匝或$\Phi$已定义为相应磁链，
则
\[
\Phi_{12}=M_{12}I,\qquad
\Phi_{21}=M_{21}I,
\]
故
\[
\boxed{\Phi_{12}=\Phi_{21}}.
\]
\crosscheck 多匝回路需明确比较的是总\emph{磁链}
而非随意某一匝的单匝磁通量。

\sourceproblem{9.2.13}{238}{238}
\problemtext 某一线圈电流在$0.002$ s内均匀地由$10$ A增至$12$ A，
其间自感电动势大小为400 V。
求该线圈的自感系数$L$。
\workings
线圈自感电动势
$|\mathcal E_L|=L|\Delta I/\Delta t|$，
故
\[
\boxed{L=\frac{400}{(12-10)/0.002}
=\frac{400}{1000}=0.40\,\mathrm H}.
\]
\crosscheck $[\mathrm V/(\mathrm{A/s})]=\mathrm H$；
电流增大时有向自感电动势取阻碍电流增长的方向。

\sourceproblem{9.2.14}{239}{239}
\problemtext 在真空中长直螺线管通电$I_1$时储磁能$W_1$。
若插入相对磁导率$\mu_r=4$的线性磁介质，
电流增为$I_2=2I_1$且几何不变，求$W_1:W_2$。
\workings
固定线圈匝数与几何，
线性自感$L\propto\mu$，
磁能$W=\tfrac12LI^2$，故
\[
\frac{W_2}{W_1}
=\frac{L_2}{L_1}\left(\frac{I_2}{I_1}\right)^2
=4\cdot2^2=16.
\]
\[
\boxed{W_1:W_2=1:16}.
\]
\crosscheck 此处假定磁介质的$\mu_r=4$为\emph{线性常数}；
对有磁滞的铁磁体并不一般满足能量恰按$\mu I^2$变化。
